\chapter{Il modulo \code{offload}}
\label{cap:13-offload}

Il modulo \code{offload} consente a un'istanza di \nf{} di inoltrare
un'invocazione sincrona a un'altra istanza remota anzich\'e eseguirla localmente.
Lo scenario di riferimento \`e edge$\to$cloud: un nodo periferico con capacit\`a
limitata che, quando non riesce ad ammettere una richiesta, la delega a
un'istanza pi\`u capiente~\cite{satya2017edge, shi2016edge}.

\`E fra i moduli pi\`u piccoli che facciano qualcosa di sostanziale: 470 righe in
tre classi (configurazione, propriet\`a, gateway), di cui 411 nel gateway; una
parte consistente di queste \`e un percorso di compatibilit\`a per i consumatori
che forniscono un registro di metriche anzich\'e il contratto del nucleo. La
brevit\`a \`e conseguenza diretta di due rinunce, entrambe deliberate, che il
capitolo argomenta, e del fatto che la decisione di scaricare non vive nel
modulo: il modulo implementa \code{OffloadGateway} (un'interfaccia del contratto
del nucleo, con un'implementazione nulla di default) e fa la chiamata remota; chi
decide \`e il nucleo.

\begin{figure}[htbp]
  \centering
  \resizebox{0.95\textwidth}{!}{% nanofaas offload: an edge control plane transparently proxies sync invocations
% to a cloud control plane under pressure (or by eager per-function policy).
% Background-layer containers are declared AFTER all blur-shadow nodes (see the
% ordering note in fig-architecture.tex).
\documentclass[border=8pt]{standalone}
\usepackage{nanofaas-figs}
\begin{document}
\begin{tikzpicture}[nf, node distance=6mm]

  % ---------- client ----------
  \node[client, minimum width=20mm, minimum height=16mm] (client) {Client};

  % ---------- edge control plane (inner nodes) ----------
  \node[edgebox, right=24mm of client, minimum width=34mm] (sq)
    {sync-queue admission\\{\scriptsize DEPTH $\cdot$ EST\_WAIT}};
  \node[edgebox, below=13mm of sq, minimum width=34mm] (epods)
    {warm pods\\{\scriptsize local dispatch}};
  \node[edgebox, right=13mm of sq, minimum width=24mm] (proxy)
    {offload\\proxy};

  % ---------- cloud control plane (inner nodes) ----------
  \node[cloudbox, right=54mm of proxy, minimum width=30mm] (cpods)
    {warm pods\\{\scriptsize cloud dispatch}};

  % ---------- failure semantics (shadowed; before the containers) ----------
  \node[danger, above=8mm of proxy, minimum width=40mm, font=\sffamily\scriptsize] (fail)
    {no local fallback\\remote failure $\Rightarrow$ \texttt{502 / 504}};

  % ---------- trigger legend (shadowed; before the containers) ----------
  \node[neutral, anchor=north west, align=left, font=\sffamily\scriptsize,
        minimum height=0mm, text width=88mm] (leg)
    at ($(epods.south west)+(-2mm,-15mm)$)
    {\textbf{Offload trigger} (single hop, no fan-out):\\
     \;$\bullet$\; \texttt{EAGER} — per-function policy (\texttt{offload.enabled=true})\\
     \;$\bullet$\; \texttt{DEPTH} / \texttt{EST\_WAIT} — sync-queue backpressure rejection\\
     Counters: \texttt{nanofaas\_offload\_total\{function,trigger\}},
     \texttt{nanofaas\_offload\_failure\_total}};

  % ---------- grouping containers (background layer, AFTER all shadowed nodes) ----------
  \begin{scope}[on background layer]
    \node[container, draw=nfblue, line width=1.5pt, fill=nfblue!9, fit=(sq)(epods)(proxy)] (edge) {};
    \node[container, draw=nfgreen, line width=1.5pt, fill=nfgreen!12, fit=(cpods)] (cloud) {};
  \end{scope}
  \node[clabel, text=nfblue, anchor=south west] at ($(edge.north west)+(1mm,0.8mm)$)
    {Edge control plane};
  \node[clabel, text=nfgreen, anchor=south west] at ($(cloud.north west)+(1mm,0.8mm)$)
    {Cloud control plane};

  % ---------- request flow ----------
  \draw[flow] (client.east) -- (client.east-|edge.west)
    node[tag, midway, above] {\texttt{:invoke}};
  \draw[flow] (sq.south) -- (epods.north)
    node[tag, midway, right=0.5mm, fill=white, inner sep=1pt] {admit};
  \draw[flow] (sq.east) -- (proxy.west)
    node[tag, midway, above, align=center] {reject /\\eager};

  % the offload hop, highlighted
  \draw[offloadarr] (proxy.east) -- (proxy.east-|cloud.west)
    node[text=nfamber, font=\sffamily\scriptsize, midway, above=0.6mm, align=center]
      {transparent proxy $\cdot$ single hop\\\texttt{X-NanoFaaS-Offloaded}};

  % ---------- responses ----------
  \draw[ret] (cpods.south) to[out=-90,in=-90,looseness=0.55]
    node[tagc, pos=0.5, below] {response — transparent to the client} (client.south);

  % ---------- failure link ----------
  \draw[flowc, draw=nfred, dashed] (fail.south) -- (proxy.north);

\end{tikzpicture}
\end{document}
}
  \caption{Offload trasparente edge$\to$cloud: attivazione eager per politica di
  funzione oppure reattiva su rifiuto della coda sincrona; salto singolo, nessun
  fallback locale.}
  \label{fig:offload}
\end{figure}

\section{Quando si scarica}

L'attivazione \`e decisa in \code{ReactiveInvocationCoordinator}, nel nucleo, sul
percorso di ammissione sincrona (\cref{cap:04-architettura-control-plane}), e ha
due innesti. In entrambi i casi vale una precondizione, valutata per prima: la
richiesta non deve provenire a sua volta da un salto di offload, e il gateway deve
essere attivo.

\subsection{Eager: la politica della funzione}

\begin{lstlisting}[language=Java]
@Override
public boolean shouldOffloadEagerly(FunctionSpec spec) {
    OffloadPolicy policy = spec.offload();
    return policy != null
            && !Boolean.FALSE.equals(policy.enabled())
            && OffloadPolicy.MODE_ALWAYS.equalsIgnoreCase(policy.mode())
            && targetUrl(spec) != null;
}
\end{lstlisting}

Una funzione con \code{offload.mode = "always"} viene sempre inoltrata, senza
tentare l'esecuzione locale. \`E la dichiarazione, da parte dell'operatore, che
quella funzione non ha senso al bordo --- per esempio perch\'e richiede risorse
che il nodo periferico non ha.

\subsection{Pressure: il rifiuto come segnale}

\begin{lstlisting}[language=Java]
try {
    admitLocally(executionRecord);
} catch (SyncQueueRejectedException ex) {
    OffloadTrigger trigger = pressureTrigger(ex.reason());
    if (offloadable && trigger != null && offloadGateway.shouldOffloadOnPressure(spec)) {
        startOffload(executionRecord, spec, trigger, context, offloadedTarget);
        return;
    }
    throw ex;
}

private static OffloadTrigger pressureTrigger(SyncQueueRejectReason reason) {
    return switch (reason) {
        case DEPTH -> OffloadTrigger.DEPTH;
        case EST_WAIT -> OffloadTrigger.EST_WAIT;
        case TIMEOUT -> null;
    };
}
\end{lstlisting}

Questa \`e la parte concettualmente interessante. La decisione di scaricare
\emph{non} \`e presa da un modello di costo che confronti la latenza locale
attesa con quella remota, come nella formulazione classica del problema di
offloading. \`E presa dal \emph{rifiuto} della coda di ammissione locale.

L'eccezione \code{SyncQueueRejectedException} \`e lanciata \emph{soltanto} dal
gateway della coda sincrona. Ne segue una conseguenza pratica che il nome del
modulo non lascia intuire: l'offload reattivo funziona solo quando l'ammissione
sincrona passa per la coda sincrona, cio\`e nel profilo \code{SYNC\_QUEUE}
(\cref{sec:09-profili}). Nei profili \code{FUNCTION\_QUEUE} e \code{DIRECT} una
funzione che trova coda o slot esauriti riceve il \code{429} del nucleo e nulla
viene scaricato: resta operativo il solo offload eager. Con entrambi i moduli di
coda nell'artefatto il profilo derivato \`e \code{FUNCTION\_QUEUE}, quindi
l'offload reattivo richiede in quel caso di dichiarare esplicitamente il profilo.

Il vantaggio \`e che non richiede alcuna calibrazione preventiva: la coda
sincrona sa gi\`a stimare l'attesa e sa gi\`a decidere che \`e eccessiva
(\cref{cap:10-sync-queue}); l'offload riusa quella decisione anzich\'e
ripeterla con parametri propri. Il costo \`e che la decisione \`e binaria e
tardiva: non c'\`e modo di scaricare \emph{preventivamente} una frazione del
traffico per evitare la congestione, e non c'\`e alcuna stima della latenza
remota nella decisione.

\subsection{Perch\'e \code{TIMEOUT} non \`e pressione}

Il \code{switch} mappa esplicitamente \code{TIMEOUT} a \code{null}. Le prime due
ragioni di rifiuto avvengono \emph{prima} che la richiesta occupi la coda: la
piattaforma dichiara di non poterla accettare, e inoltrarla altrove \`e una
risposta sensata. \code{TIMEOUT} avviene \emph{dopo} che la richiesta ha
occupato la coda per l'intera durata concessa: il budget temporale \`e gi\`a
esaurito, e iniziare adesso una chiamata remota consumerebbe altro tempo per
consegnare comunque una risposta tardiva.

\section{Le due rinunce}

\subsection{Salto singolo}

Ogni richiesta inoltrata porta l'header \code{X-NanoFaaS-Offload-Hop: 1}, e il
control plane ricevente lo interpreta come divieto di ulteriore inoltro:

\begin{lstlisting}[language=Java]
boolean offloadable = !context.offloadedHop() && offloadGateway.enabled();
\end{lstlisting}

La ragione \`e la terminazione. Con $n$ istanze configurate in catena o in ciclo,
una richiesta potrebbe rimbalzare indefinitamente, e ogni rimbalzo consuma parte
del budget temporale del chiamante originale senza avvicinarsi a un'esecuzione.
Un contatore di salti con limite $k > 1$ sarebbe pi\`u generale e richiederebbe
di decidere $k$; il valore 1 \`e l'unica scelta che non richiede una decisione.

\subsection{Nessun fallback locale}

Se la chiamata remota fallisce, l'esecuzione fallisce. Il chiamante riceve
\code{502 Bad Gateway} o \code{504 Gateway Timeout}. Non viene ritentata
localmente.

Il commento nel codice del nucleo motiva la scelta sul piano meccanico:

\begin{lstlisting}[language=Java]
// Bypasses the local queue entirely: no local concurrency slots are consumed,
// so completion goes through the offload-specific path (no slot release, no retry).
\end{lstlisting}

Una richiesta inoltrata non ha mai acquisito uno slot locale. Un fallback la
farebbe rientrare nella coda locale --- la stessa che l'aveva rifiutata, e che
nel frattempo non si \`e svuotata --- riproducendo il rifiuto con il budget
temporale ormai consumato. Il fallback trasformerebbe un \code{502} rapido in un
\code{429} tardivo, che \`e strettamente peggiore.

C'\`e anche una ragione di onest\`a diagnostica: con il fallback, un'istanza
remota irraggiungibile si manifesterebbe come un aumento di latenza locale
anzich\'e come un errore, e la configurazione errata resterebbe invisibile.

\section{La semantica degli errori}

La distinzione fra ``la funzione remota ha fallito'' e ``l'inoltro ha fallito''
\`e centrale, ed \`e implementata con cura:

\begin{lstlisting}[language=Java]
// A marked response is the function's own answer, whatever its status —
// read it down the same body-parsing path as a plain 2xx, before any
// status-based branching (a marker-bearing 404 is a function decision,
// not "unregistered function").
if (functionDecided || response.statusCode().is2xxSuccessful()) {
    return response.bodyToMono(InvocationResponse.class)
            .map(this::toResult)
            .switchIfEmpty(Mono.error(new OffloadFailedException(target, false,
                    "empty response body from remote " + target)));
}
if (response.statusCode().value() == 404) {
    return response.releaseBody().then(Mono.error(new OffloadFailedException(target, false,
            "function '" + task.functionName() + "' not registered on remote " + target)));
}
\end{lstlisting}

Il caso critico \`e il \code{404}. Un \code{404} \emph{senza} marcatore significa
che la funzione non \`e registrata sull'istanza remota: \`e un errore di
configurazione della piattaforma, e va segnalato come fallimento di offload
(\code{502}), non propagato al chiamante come se la sua richiesta fosse
malformata. Un \code{404} \emph{con} marcatore \code{X-NanoFaaS-Function-Status}
\`e invece la risposta deliberata della funzione remota, e va propagato tale e
quale.

L'ordine dei rami \`e quindi essenziale: il marcatore viene consultato prima di
qualunque diramazione sul codice di stato.

\subsection{Che cosa attraversa il salto}

Il gateway inoltra la richiesta originale come envelope JSON verso
\code{/v1/functions/\{nome\}:invoke} dell'istanza remota, e vi aggiunge il
marcatore di salto e, se presenti, gli header di tracciamento
(\code{X-Trace-Id}, \code{traceparent}, \code{tracestate}). Il control plane
ricevente ricostruisce gli header applicativi della richiesta dal trasporto HTTP
del salto, quindi il gateway copia gli header applicativi del chiamante come veri
header HTTP: senza questa copia la funzione remota li perderebbe.

La copia \`e difensiva. Sono esclusi: l'inquadramento del trasporto che il client
HTTP possiede (\code{host}, \code{content-length}, \code{transfer-encoding} e
soprattutto \code{content-type}, che deve continuare a descrivere l'envelope JSON
e non il corpo originale del chiamante, e con esso la negoziazione e la codifica
del contenuto), gli header che il gateway imposta da s\'e o che il ricevente
lega a parametri dedicati (identificativi di esecuzione e di tentativo, timeout,
chiave di idempotenza, marcatore di salto, tracciamento), gli header
\emph{hop-by-hop} definiti da RFC~9110 e ogni header \code{proxy-*}. L'elenco
rispecchia di proposito quello del controller ricevente: inoltrare un nome che il
ricevente scarterebbe non serve, e un valore applicativo non deve poter fingersi
un header di controllo.

\subsection{Il margine di timeout}

\begin{lstlisting}[language=Java]
// ponytail: fixed margin so the gateway's remote budget expires before the
// coordinator's local wait, making 504 (not a local "timeout") deterministic
private static final long TIMEOUT_MARGIN_MS = 50;
...
long timeoutMs = Math.max(1, timeoutBudgetMs - TIMEOUT_MARGIN_MS);
\end{lstlisting}

Due timeout corrono in parallelo: quello locale del coordinatore, che attende il
futuro di completamento, e quello della chiamata remota. Se scadessero
contemporaneamente, l'esito dipenderebbe da una corsa e il chiamante riceverebbe
talvolta \code{504} e talvolta una risposta di timeout locale. Il margine fisso
di 50\,ms garantisce che il timeout remoto scatti per primo, e quindi che la
diagnosi sia deterministica.

Il margine ha una condizione che il codice non enuncia: al gateway viene passato
come budget il timeout della \emph{specifica} della funzione, mentre l'attesa
locale usa l'eventuale timeout indicato dal chiamante per quella richiesta. Se il
chiamante chiede un timeout inferiore a quello della specifica, l'attesa locale
scade prima e il margine non ordina pi\`u i due eventi; la determinazione del
\code{504} vale, in quel caso, solo per l'attesa di specifica.
\begin{damisurare}
\todomis{esito osservato per una richiesta offload con timeout per-richiesta inferiore a quello della specifica}
\end{damisurare}

\`E un dettaglio piccolo con una propriet\`a importante: rende \emph{riproducibile}
un comportamento in condizioni di errore, che \`e la classe di comportamenti pi\`u
difficile da testare e la pi\`u facile da rendere non deterministica per
distrazione.

\subsection{Chi possiede l'input}

Un'esecuzione scaricata non acquisisce alcun \code{DispatchLease}, ma tiene comunque
una risorsa: l'input fisico della richiesta. La chiamata remota ne \`e proprietaria
fino al proprio segnale terminale, non fino a quando la vista locale \`e ancora
annullabile. Una cancellazione amministrativa conclude soltanto la vista locale
dell'esecuzione; il publisher della chiamata remota pu\`o ignorare la cancellazione
e conservare la richiesta, quindi l'input viene chiuso solo quando la chiamata
termina, in qualunque modo. \`E la stessa regola di propriet\`a che il nucleo applica
agli altri percorsi (\cref{cap:04-architettura-control-plane}).

L'esito di un'esecuzione scaricata \`e definitivo: il completamento passa per un
percorso dedicato del nucleo che non ritenta e non rilascia slot, e un fallimento
dell'inoltro conclude l'esecuzione con il codice \code{OFFLOAD\_FAILED} o, se
scaduto il budget, \code{OFFLOAD\_TIMEOUT}. Il chiamante li vede come \code{502} e
\code{504}.

\section{Osservabilit\`a}

\begin{center}
\small
\begin{adjustbox}{max width=\textwidth}
\begin{tabular}{@{}lll@{}}
\toprule
\textbf{Segnale} & \textbf{Tipo} & \textbf{Etichette} \\
\midrule
\code{nanofaas.offload} & counter & \code{function}, \code{trigger}
(\code{eager} / \code{depth} / \code{est\_wait}) \\
\code{nanofaas.offload.failure} & counter & \code{function} \\
\code{X-NanoFaaS-Offloaded} & header & URL dell'istanza remota, anche sulle
risposte \code{502} e \code{504} \\
\bottomrule
\end{tabular}
\end{adjustbox}
\end{center}

L'header di risposta \`e ci\`o che rende l'offload \emph{trasparente ma non
invisibile}: il chiamante non deve fare nulla di diverso, ma se lo desidera pu\`o
sapere che la sua richiesta \`e stata servita altrove e da chi. \`E la stessa
distinzione fra trasparenza e opacit\`a che il \cref{cap:07-modalita-esecuzione}
applica alla degradazione da \code{DEPLOYMENT} a \code{EXTERNAL}.

I due contatori sono legati alla \emph{generazione} della funzione attiva
all'inizio dell'offload. Il tentativo \`e contato alla sottoscrizione della
chiamata e non quando la pipeline viene costruita, e un fallimento \`e contato solo
per un tentativo gi\`a contato; un tentativo appartenente a una generazione ritirata
scarta silenziosamente i propri aggiornamenti, cos\`i che una funzione cancellata
non veda risorgere la serie (invariante I7 del ciclo di vita).

L'etichetta \code{trigger} sul contatore consente di distinguere le tre cause di
attivazione, il che \`e necessario per interpretare il dato: un'alta frequenza di
\code{eager} \`e una politica che funziona, un'alta frequenza di \code{depth} \`e
un nodo periferico sotto-dimensionato.

\section{Configurazione}

\begin{center}
\small
\begin{adjustbox}{max width=\textwidth}
\begin{tabular}{@{}llp{6.6cm}@{}}
\toprule
\textbf{Chiave} & \textbf{Default} & \textbf{Effetto} \\
\midrule
\code{nanofaas.offload.enabled} & \code{true} & Interruttore generale; il gateway
resta comunque inerte senza un target. \\
\code{nanofaas.offload.target-url} & --- & URL base dell'istanza remota. \\
\code{nanofaas.offload.pressure-enabled} & \code{true} & Attiva l'offload
reattivo; quello eager per funzione resta comunque. \\
\midrule
\multicolumn{3}{@{}l@{}}{\emph{Per funzione, nel blocco \code{offload}:}} \\
\code{enabled} & --- & \code{false} esclude la funzione. \\
\code{targetUrl} & --- & Sovrascrive il target globale; la sua sola presenza
attiva l'offload per quella funzione. \\
\code{mode} & \code{pressure} & \code{always} per l'offload eager. \\
\bottomrule
\end{tabular}
\end{adjustbox}
\end{center}

Il default \code{enabled = true} in assenza di un target \`e apparentemente
incoerente ed \`e invece corretto: il modulo \`e attivo ma non ha dove inoltrare,
quindi non fa nulla. Configurare un target \`e l'unico atto necessario, e non
occorre ricordarsi anche di attivare l'interruttore. Il commento nel codice lo
enuncia: \emph{``even when true the gateway stays inert until a target URL is
configured''}.

\section{Ci\`o che il modulo non fa}

\begin{itemize}[leftmargin=*]
  \item \textbf{Non modella il costo.} La formulazione classica del problema di
  offloading valuta se il trasferimento convenga, dati il costo di rete e il
  guadagno di esecuzione. Qui la decisione \`e binaria e derivata dal rifiuto
  locale.

  \item \textbf{Non sceglie fra pi\`u destinazioni.} Il target \`e uno, globale
  o per funzione. Non esiste bilanciamento fra istanze remote.

  \item \textbf{Non scarica il percorso asincrono.} Solo \code{:invoke} \`e
  soggetto a offload. L'invocazione asincrona ha gi\`a una coda che assorbe la
  pressione, e il chiamante non sta attendendo.
\end{itemize}

\begin{damisurare}
Il beneficio dell'offload sotto pressione non \`e stato caratterizzato
sperimentalmente: non esiste una campagna che confronti latenza end-to-end e
tasso di rifiuto di un nodo periferico con e senza offload attivo, al variare
della latenza di rete verso il target. \`E l'esperimento naturale per questo
modulo e non \`e stato eseguito.
\todomis{p95 end-to-end e tasso di rifiuto con e senza offload, al variare del
RTT verso l'istanza remota}
\end{damisurare}
